% Draft addition to Section 7.6 of relay/paper/main.tex, to follow the existing Gaussian
% paragraph. Every number below comes from paper/experiments/exp6_robustness.py at 25 seeds
% per cell, 72 cells per configuration. Raw files: exp6_robustness_satcap{15,30}.csv,
% exp6_robustness_undercount0.{5,6,7,8,9}.csv, exp6_robustness_undercount0.5_presence.csv.
%
% NIRAJ: these are new results. Read them, re-run them yourself, and decide whether the
% reading below is one you are willing to defend. They are your claims, not mine.

That test is the easy one, and reported alone it would flatter the controller. A zero-mean
error averages out over the smoothing window by construction, whereas
Section~\ref{sec:perception} measures a \emph{biased} one: on motorcycle-dense footage the
detector reports 26 vehicles in a scene holding several hundred, and sampled counts are still
climbing 83\% between 640~px and 1600~px without converging. A biased estimator does not
average out. We therefore repeat the full grid under count models that are wrong in the
direction the detector is actually wrong. In every case the queues the controller is managing
stay exact, so what is measured is the control response to a wrong observation.

The first model saturates: the controller observes $\hat q = q/(1 + q/C)$, near-exact on a
lightly loaded approach and increasingly short as the approach fills, which is the shape of the
measured failure. At both $C = 15$ and $C = 30$~vehicles the result is unchanged, 71 of 72
cells, a worst margin against the equal-split plan of $+26.6\%$, and a mean cost across the grid
of $-0.3$ percentage points. The largest loss in any single cell is 1.6 points. The one cell the
controller loses in fact improves, from $-26.6\%$ to $-4.2\%$, because compressing the dominant
approach's count stops the policy contending so hard against the fairness bound.

The second model is a flat multiplicative loss, $\hat q = \lfloor \kappa q \rceil$ applied
equally to every approach. Down to $\kappa = 0.6$ it costs nothing measurable: 71 of 72 cells at
$\kappa = 0.9$, $0.8$, $0.7$ and $0.6$, with a mean cost of at most $0.1$ percentage points.
This is a property of the formulation rather than a lucky range. The degree of saturation of
Section~\ref{sec:floor} is a ratio of an arrival estimate to a discharge estimate, both derived
from the same count sequence, so a uniform scale factor cancels before it reaches a decision.

At $\kappa = 0.5$ the controller collapses, to 50 of 72 cells and a mean cost of 32.7 points,
falling behind even the equal-split plan in the worst cell ($-11.1\%$). The discontinuity
between $0.6$ and $0.5$ is not a magnitude effect. At $\kappa = 0.5$ a lone vehicle rounds to
zero, so an occupied approach reports empty and empty-phase skipping declines to serve it; at
$\kappa = 0.6$ it survives rounding. Repeating $\kappa = 0.5$ with the single change that an
occupied approach may never report empty restores the result completely: 71 of 72 cells, worst
margin $+26.4\%$, mean cost $0.2$ points. The 50\% magnitude error was never the problem.

We state the resulting requirement precisely, because it is narrower and more testable than
"the detector must be accurate". What this controller needs from perception is an unbiased
\emph{presence} signal at short queue lengths, not an accurate absolute count. A detector that
undercounts a crowd by an order of magnitude is tolerable, because both the pressure ordering
and the saturation ratio survive it; a detector that cannot distinguish an empty approach from
one holding a single vehicle is not, whatever its aggregate accuracy. That is a condition a
field trial can check on a candidate detector before any signal is connected to it, and it is a
condition the perception stack of Section~\ref{sec:perception} has not yet been shown to meet.
